\honest{Religion's Claim to be Non-Disprovable}{Eliezer Yudkowsky}{2007}

The earliest account I know of a scientific experiment is Elijah's contest with the priests of Baal: each side lays out a bull on the same kind of wood, and the true God is to send fire. The priests of Baal are the control group, and Elijah soaks his own altar with water, accepting the burden of proof. The fire falls on Elijah's altar, the people shout ``The Lord is God!'', and the 450 priests are killed. Joking, I call that ``stern, but necessary'': otherwise they would have started saying that religion is ``a separate magisterium which can be neither proven nor disproven.'' \nb{The phrase is Stephen Jay Gould's. In 1997 he argued that science and religion each have ``a legitimate magisterium, or domain of teaching authority,'' and that these do not overlap. The post never names him.}

In the old days people believed their religions as fact. The archaeologists who searched for Noah's Ark expected they might become famous. Only when the evidence went against them did believers make what William Bartley called ``the retreat to commitment,'' which I gloss as ``I believe because I believe.'' \nb{Bartley's phrase named an argument: that every position, rationalism included, rests on some arbitrary starting commitment. The post gives the conclusion without the argument.}

The Old Testament is a ``culture dump'' that includes models of the universe. Most religions in history tell of events that would be ``completely unmistakable evidence'' if they had happened, and separating religion from fact is ``a recent and strictly Western concept.'' Orthodox Judaism holds that God spoke at Mount Sinai in a thundering voice; if that happened, it would be plain evidence of some superhuman being, though not that it was God, or good. \nb{Every religious example in the post comes from the Bible.}

With no libraries to check it, the Old Testament could claim the Exodus, of which Egyptologists found no trace. They did find that Egypt ruled much of Canaan at the supposed time, a large error that no one could then catch. \nb{Modern scholarship, as Wikipedia summarizes it, broadly agrees.} The Romans kept records, so the New Testament claims only smaller miracles, like casting out the spirit that made a boy froth at the mouth. That is still ordinary evidence: if epilepsy is caused by demons, and demons flee a true prophet but not a fraud, then the end of the fit shows Jesus is a true prophet. \nb{The New Testament also reports a darkness ``over all the earth'' at the crucifixion, and dead saints rising and appearing ``unto many'' in Jerusalem.}

Religion once made claims about everything: law before legislatures, history before historians, sexual morals, forms of government, and science from the classing of animals to the formation of stars. Better institutions took over each field, and ethics is what is left. But the Bible's ethics are obsolete too: it approves of slavery, and the killing of Egypt's first-born is worse than any error about grasshoppers' legs. \nb{That is a moral judgment, not evidence. It shows that the Bible's ethics can be criticized, not that religious claims can be disproved by observation, which is the post's thesis.} If you say the Earth is flat, people think you mad; if you say the Bible is your source of ethics, they do not. I add, with no evidence: ``Most people's concept of rationality is determined by what they think they can get away with,'' and everyone has agreed not to notice.

So the idea that religion cannot be proven or disproven is ``a Big Lie,'' false to history, to scripture, to what children are told and to what most believers hold. It is like a defendant who, faced with the bloody axe, says his innocence cannot be disproved by mere evidence, since it is a separate magisterium. \nb{Gould's proposal was about how the two domains should be related, and he said it was the standard attitude of the major Western religions. The post's history does not answer the first claim, and it gives no evidence for its own claim about what most believers hold.}

My footnote claims that the Old Testament shows no wonder at the universe, being busy ``laying down the death penalty for women who wore mens clothing.'' \nb{Psalms 8 and 19 and the closing chapters of Job express wonder at the heavens and at nature. Deuteronomy 22:5, which forbids cross-dressing, sets no penalty.}
